% !TeX root = ../机器学习笔记.tex
% 由主文档通过 \input 载入；请编译主文档。
\setcounter{chapter}{21} % 保留原书章号。
\chapter{图数据、推荐与排序}
\label{chap:ml-graphs}

图记录对象之间的联系，推荐利用用户与物品的关系作出预测，排序则决定展示顺序。关系数据中的样本通常不独立，信息可用范围和评价单位必须与任务一起规定。

\section{图表示与学习任务}

图 $G=(V,E)$ 的节点代表对象，边代表联系，可带方向与权重。邻接矩阵 $A_{ij}$ 记录从节点 $i$ 到节点 $j$ 的连接；无向图矩阵对称。节点分类预测节点标签，链接预测判断尚未观测的边，图级预测则为整个图赋值。相同结构可用于不同目标，不能仅看数据形式判断任务。

训练时已知测试节点的特征与允许使用的边、只隐藏其标签，属于一种传导式设定；要求适用于以后到达的新节点或新图，则通常需要归纳式规则。共享图结构不一定是泄漏，关键在于这种共享是否属于预先规定、符合应用的信息条件。

\section{随机游走与 PageRank}

将出边权重归一为行随机矩阵 $P$；无出边节点的行必须用明确的概率分布补齐。给定全分量正且和为一的列向量 $v$、$0<\alpha<1$，PageRank 满足
\begin{equation}
 p=\alpha P^\top p+(1-\alpha)v.
 \label{eq:ml-pagerank}
\end{equation}
它描述以概率 $\alpha$ 沿图移动、以概率 $1-\alpha$ 按 $v$ 重新跳转的长期访问权重。

\begin{proposition}
映射 $T(p)=\alpha P^\top p+(1-\alpha)v$ 在概率单纯形上有唯一固定点，且从任意概率初值迭代收敛到它。
\end{proposition}
\begin{proof}
行随机性保证 $P^\top$ 保持概率向量，且
$\|P^\top(p-q)\|_1\leq\sum_{ij}P_{ij}|p_i-q_i|=\|p-q\|_1$。
因此 $\|T(p)-T(q)\|_1\leq\alpha\|p-q\|_1$。相邻两次迭代的距离至多为初始相邻距离的 $\alpha^t$ 倍，几何级数可求和，故迭代为 Cauchy 序列。单纯形完备，极限仍在其中；$T$ 连续，极限就是固定点。若两个固定点不同，其距离既等于原距离又至多为原距离的 $\alpha$ 倍，产生矛盾，故固定点唯一。再次使用压缩不等式即得误差至多按 $\alpha^t$ 倍缩小。
\end{proof}
高访问权重是一种结构重要性，不等于针对任意查询的监督相关性。

\section{协同过滤与矩阵分解}

协同过滤（collaborative filtering）利用相似用户、相似物品或共享潜在因子预测偏好。对观测评分集合 $\Omega$，常用模型
$\widehat r_{ui}=\mu+b_u+c_i+p_u^\top q_i$，
通过观测项平方损失与参数二次惩罚训练。$\mu$ 是全局基准，$b_u,c_i$ 是用户和物品偏置，$p_u,q_i$ 是潜在向量。

只在 $\Omega$ 上计算拟合损失，因为未评分既可能表示未见过，也可能表示不愿评分，不能直接当作零偏好。交替固定一组因子求另一组，或作随机梯度更新，都只保证相应数值目标的局部求解。冷启动用户或新物品缺少交互，通常需要内容特征或其他明确的默认规则。

\section{隐式反馈与学习排序}

点击等隐式反馈依赖曝光：没有点击可能是未展示而非不喜欢。负样本构造因此是目标定义的一部分。逐点方法预测每对用户物品的分数，成对方法直接鼓励正项优于负项，例如
$-\log\sigma(s_{ui}-s_{uj})$；
列表式方法考虑整个候选列表的排列质量。

设位置 $j$ 的相关等级为 $r_j$，折损累计增益为
$\mathrm{DCG}@k=\sum_{j=1}^k(2^{r_j}-1)/\log_2(j+1)$，
除以相同候选集的理想排序增益得到 NDCG，理想增益为零时需约定。改变候选集、负采样或可见历史都会改变指标含义，故不同协议下的分数不宜直接比较。

\section{图表示、传播与评价边界}

图神经网络（graph neural network，GNN）通过邻域聚合更新表示。对于无向图，一种常见归一化传播为
\begin{equation}
 H^{(l+1)}=\phi(\widetilde D^{-1/2}\widetilde A\widetilde D^{-1/2}H^{(l)}W^{(l)}),
 \label{eq:ml-graph-propagation}
\end{equation}
其中 $\widetilde A=A+I$ 加入自环，$\widetilde D$ 为其度矩阵，$H$ 的行对应节点。聚合把邻居信息引入当前节点，是否有用取决于边的语义及任务所需关系。

反复平滑可能使节点表示趋同，丢失区分信息。链接预测尤其需要在构图时去除应当隐藏的目标边；否则传播已直接利用答案。评估按节点、边还是整个图划分，代表不同泛化问题，必须随结果一同说明。

\section{例题、练习与学习检查}

\begin{example}[两节点 PageRank]
两个节点相互指向，$P=\left(\begin{smallmatrix}0&1\\1&0\end{smallmatrix}\right)$，$v=(0.8,0.2)^\top$，$\alpha=1/2$。
\end{example}
\begin{solve}
$p_1=\frac12p_2+0.4$，$p_2=\frac12p_1+0.1$，故 $p=(0.6,0.4)^\top$。原始游走具有二周期，带跳转的迭代仍按压缩性收敛。
\end{solve}

习题：用残差 $\|T(p)-p\|_1$ 推导固定点误差上界 $\|T(p)-p\|_1/(1-\alpha)$；推导一个评分样本的因子梯度；解释热门物品基线的作用；设计只用过去交互预测未来推荐效果的评价协议。





